\documentclass[12pt,a4paper]{ctexart}
\usepackage{geometry}
\geometry{left=2.5cm,right=2.5cm,top=2.5cm,bottom=2.5cm}
\usepackage{listings}
\usepackage{xcolor}
\usepackage{booktabs}
\usepackage{amsmath}
\usepackage{tikz}
\usepackage{graphicx}

\lstset{
  basicstyle=\small\ttfamily,
  keywordstyle=\color{blue}\bfseries,
  commentstyle=\color{gray}\itshape,
  numbers=left,numberstyle=\tiny\color{gray},
  frame=single,rulecolor=\color{black!30},
  breaklines=true,breakatwhitespace=true,
  language=C++
}

\title{\textbf{编译器设计专题实验报告}\\[0.5em]实验五：语义分析框架实现}
\author{张凌曜 \quad 2234214014 \quad 计算机2306}
\date{}

\begin{document}
\maketitle

\section{实验目的}
在词法分析和语法分析的基础上，构建AST（抽象语法树）和符号表，做类型检查和语义错误检测（比如重复声明、未声明变量）。

\section{实验原理}
语义分析是编译的第三阶段。AST用树形结构表示程序的语法层次，符号表记录变量名、类型、作用域这些信息，类型检查就是看赋值和运算中类型对不对得上。

本实验用递归下降解析器，解析的同时建树和填符号表。

\section{实验过程}

\subsection{AST节点}
每个AST节点有类型（FuncDecl、VarDecl、Assign、If、While等）、值（变量名、运算符）、数据类型（int/float）和子节点列表。

\subsection{符号表}
符号表用作用域栈实现。每进一个代码块（函数体、if块）压一层新作用域，退出时弹出。变量查找从内到外逐层找，先找到先用，这样能处理变量遮蔽。重复声明检测只看当前作用域内有没有同名变量。

\subsection{类型推断}
二元运算的类型推断规则很简单：任一操作数是float结果就是float，否则是int。

\section{测试结果}

测试程序：

\begin{lstlisting}[language={}]
int main() {
    int x;
    float y;
    x = 5;
    y = 3.14;
    if (x < 10) {
        x = x + 1;
    }
    while (y > 0) {
        y = y - 1.0;
    }
    print x;
    return 0;
}
\end{lstlisting}

程序输出的AST（缩进表示层级）：

\begin{lstlisting}[language={},basicstyle=\footnotesize\ttfamily]
Program
  FuncDecl(main) :int
    Params
    Block
      VarDecl(x) :int
      VarDecl(y) :float
      Assign(x) :int
        ID(x)
        IntConst(5) :int
      Assign(y) :float
        ID(y)
        FloatConst(3.14) :float
      If
        BinOp(<)
          ID(x) :int
          IntConst(10) :int
        Block
          Assign(x) :int
            ID(x)
            BinOp(+) :int
              ID(x) :int
              IntConst(1) :int
      While
        BinOp(>)
          ID(y) :float
          IntConst(0) :int
        Block
          Assign(y) :float
            ID(y)
            BinOp(-) :float
              ID(y) :float
              FloatConst(1.0) :float
      Print
        ID(x) :int
      Return
        IntConst(0) :int
\end{lstlisting}

符号表在作用域0记录了函数\texttt{main}的签名\texttt{int()}。语义检查结果：未检测到语义错误，符合预期——变量声明、赋值、if/while的类型都匹配正确。

\begin{center}
\includegraphics[width=\textwidth,height=7cm,keepaspectratio]{image/1.png}
\end{center}

\begin{center}
\includegraphics[width=\textwidth,height=6cm,keepaspectratio]{image/2.png}
\end{center}

\section{实验总结}
递归下降写起来有个容易搞混的地方：同样是\texttt{int x}后面跟分号是变量声明，跟左括号就是函数声明，需要往前多看一个Token才能确定。

整体来说AST能比较清楚地反映程序结构，符号表的作用域栈处理变量遮蔽也没问题。

\appendix
\section{附录：代码仓库}
本实验的全部源码、报告及Web演示平台已托管至GitHub：

\texttt{https://github.com/yaossss123/compilation-web}

\end{document}
